%============================================================ 6. 결과
\section{결과}
\label{sec:results}

\S\ref{sec:experiment}의 설계가 답하려는 질문은 하나다 --- 두 계층의 기여를 분리했을 때 계층별로 무엇을 얼마나 사는가 --- 기동 계층은 그물과 선회량을, 지휘관 계층은 능력 임계 위에서만 포획률을. 아래는 그 답을 세 묶음으로 보인다: 기동 계층이 사는 것(\S\ref{sec:res_g1}), 지휘관 계층이 사는 것과 그 조건(\S\ref{sec:res_g2}), 그리고 둘이 더해지는가(\S\ref{sec:res_g3}). 마지막으로 배치 관점의 지연을 보고한다(\S\ref{sec:res_latency}).

\subsection{기동 계층이 사는 것: 같은 방어를 더 적은 자원으로}
\label{sec:res_g1}

\subsubsection{기동 계층: 휴리스틱 대 U-Net 점수맵 (LLM 미사용)}
\label{sec:res_maneuver}

먼저 언어모델을 배제하고, \S\ref{sec:heuristic}의 기하 휴리스틱과 \S\ref{sec:policy}의 점수맵 정책만 비교한다. 두 조건 모두 배정은 휴리스틱이므로 차이는 오직 기동 방식이다
(Table~\ref{tab:maneuver}). 3개 포메이션 통합 기준으로 포획률은 $0.844 \to 0.886$
($+0.041$, $\dz = 0.37$)로 개선되었고 $95\,\%$ 신뢰구간이 0을 포함하지 않는다.
\textbf{다만 이 개선은 아슬아슬하다} --- 지표 10개에 대한 다중비교 보정을 걸면 구간이
$[-0.013, +0.100]$ 로 0을 포함한다(\S\ref{sec:stats}). 이 절이 보이려는 것은 포획률의
증가가 아니라, 뒤에서 보이듯 \emph{같은 정도로 막으면서 자원과 기동을 크게 줄인다}는
쪽이다.

% 자동 생성 — tools/make_paper_tables.py. 손으로 고치지 말 것.
\begin{table*}[tp]
\centering
\caption{기동 계층(U-Net 점수맵 + GRPO)의 단독 기여. 두 조건 모두 배정은 휴리스틱이며 \textbf{LLM 을 전혀 사용하지 않는다} --- 따라서 이 결과는 외부 API 에 의존하지 않고 재현 가능하다. 파상 포메이션에서 가장 크다.}
\label{tab:maneuver}
\footnotesize
\resizebox{\ifdim\width>\linewidth\linewidth\else\width\fi}{!}{%
\begin{tabular}{lcccc}
\toprule
포메이션 & 휴리스틱 기동 & U-Net 기동 & $\Delta$ [95\,\% CI] & $d_z$ \\
\midrule
집중 & 0.997 & 0.997 & +0.000 [-0.010, +0.010] & +0.00 \\
양동 & 0.787 & 0.753 & -0.033 [-0.077, +0.000] & -0.31 \\
파상 & 0.750 & 0.907 & +0.157 [+0.043, +0.270]$^{*}$ & +0.48 \\
\textbf{전체} & 0.844 & 0.886 & +0.041 [+0.003, +0.082]$^{*}$ & +0.37 \\
\bottomrule
\end{tabular}
}
\\[3pt]
\begin{minipage}{\linewidth}\footnotesize
$^{*}$ 신뢰구간이 0 을 포함하지 않음. 전체 행은 3개 포메이션 통합.
\end{minipage}
\end{table*}


\begin{figure*}[tp]
\centering
\includegraphics[width=0.9\textwidth]{fig14.png}
\caption{휴리스틱 대 U-Net 기동의 지표 프로파일(바깥일수록 우수): (a) 전체, (b) 집중,
         (c) 양동, (d) 파상}
\label{fig:radar_maneuver}
\end{figure*}

같은 비교를 10개 지표 전부에 걸쳐 한눈에 보인 것이 Fig.~\ref{fig:radar_maneuver}다. 각
축은 지표별 조건 평균을 평가한 8개 조건의 범위로 0--1 정규화한 값이며(낮을수록 좋은 지표는
방향을 뒤집어 1 = 8개 조건 중 최선), 두 다각형이 겹치는 축은 기동 방식이 결과를 바꾸지
않는 지표다. 전체(a)에서 U-Net 다각형은 포획률·돌파·포획 이격거리·포획 시각·포획당
그물·그물 걸림·총 선회량 축에서 바깥에 있고, 충돌률·전개 시 적 거리 축에서는 거의 겹치며,
포획당 이동거리 축에서만 안쪽이다 --- 자리를 골라 그물벽을 놓는 대신 더 멀리 움직이는 셈이다.

포메이션별로 보면 기여가 고르지 않다. 파상 포메이션에서 $+0.157$($\dz = 0.48$)로 가장 크고,
집중 포메이션에서는 $0.000$, 양동 포메이션에서는 $-0.033$ 으로 오히려 소폭 낮다(다만 후자의
신뢰구간은 0을 포함하므로 악화라고 말할 수 없다).

이 분포는 포메이션의 구조와 대응한다. 파상 공격은 이웃한 방위의 3개 단이 거리차를 두고
도달하므로, 앞 단을 막은 그물벽의 \emph{위치}가 뒤 단까지 막을 수 있는지를 좌우한다.
해면 전체를 평가하여 자리를 고르는 점수맵 방식이 여기서 직접적으로 값을 한다. 반면 양동
공격은 세 방위에서 동시에 도래하여 방어정 수와 무리 수가 같으므로, 결정적인 것은 어느
방어정을 어느 무리에 붙이는가라는 \emph{배정} 판단이다. 배정이 어긋나면 기동 계층이 아무리
좋은 자리를 골라도 애초에 잘못된 무리를 막게 되므로 보정되지 않는다. 집중 포메이션은
뒤에서 보이듯 천장 구간이다.

\textbf{운용 함의.} 언어모델 없이 기동 계층만 교체해도 같은 방어 성과를 더 적은 그물과 기동으로 얻는다. 이 조건은 외부 API가 없는 온보드 구성에서 먼저 적용할 수 있는 부분이다.

\subsection{지휘관 계층이 사는 것과 그 조건}
\label{sec:res_g2}

\subsubsection{전체 조건 비교}
\label{sec:res_main}

Table~\ref{tab:main}과 Fig.~\ref{fig:condbars}에 8개 조건의 포획률을 정리하였다. 그림의
무채색 두 막대가 언어모델 없는 기준선이고 그 오른쪽은 모델 규모 오름차순이며, 오차막대는
시드 부트스트랩 95\,\% 신뢰구간이다.

% 자동 생성 — tools/make_paper_tables.py. 손으로 고치지 말 것.
\begin{table*}[tp]
\centering
\caption{조건별 포획률. 평균과 부트스트랩 95\,\% 신뢰구간, baseline 대비 대응표본 차이. 3개 포메이션 통합(각 조건 30시드).}
\label{tab:main}
\footnotesize
\setlength{\tabcolsep}{4pt}%
\resizebox{\ifdim\width>\linewidth\linewidth\else\width\fi}{!}{%
\begin{tabular}{llccc}
\toprule
지휘관(배정) & 기동 & 포획률 [95\,\% CI] & $\Delta$ vs baseline & $d_z$ \\
\midrule
휴리스틱 (baseline) & 휴리스틱 & 0.844 [0.808, 0.880] & --- & --- \\
휴리스틱 & U-Net 점수맵 & 0.886 [0.847, 0.918] & +0.041 [+0.003, +0.082]$^{*}$ & +0.37 \\
Qwen2.5 7B & 휴리스틱 & 0.652 [0.571, 0.716] & -0.192 [-0.264, -0.129]$^{*}$ & -1.00 \\
Qwen2.5 14B & 휴리스틱 & 0.811 [0.766, 0.854] & -0.033 [-0.086, +0.017] & -0.23 \\
Gemini 3.5 Flash-Lite & 휴리스틱 & 0.904 [0.864, 0.937] & +0.060 [+0.024, +0.100]$^{*}$ & +0.56 \\
Qwen2.5 7B & U-Net 점수맵 & 0.641 [0.582, 0.693] & -0.203 [-0.274, -0.138]$^{*}$ & -1.03 \\
Qwen2.5 14B & U-Net 점수맵 & 0.754 [0.698, 0.806] & -0.090 [-0.153, -0.033]$^{*}$ & -0.54 \\
Gemini 3.5 Flash-Lite & U-Net 점수맵 & 0.907 [0.869, 0.937] & +0.062 [+0.013, +0.113]$^{*}$ & +0.44 \\
\bottomrule
\end{tabular}
}
\\[3pt]
\begin{minipage}{\linewidth}\footnotesize
$^{*}$ 신뢰구간이 0 을 포함하지 않음. $d_z$ 는 대응표본 효과크기. 첫 두 열이 $2\times2$ 설계의 두 축이다 --- 같은 지휘관 행 두 개를 비교하면 기동 계층의 기여가, 같은 기동 열을 비교하면 지휘관의 기여가 나온다. Qwen2.5 는 로컬(단일 RTX 4070) 서빙, Gemini 는 클라우드 API 다.
\end{minipage}
\end{table*}


\begin{figure*}[tbp]
\centering
\includegraphics[width=0.9\textwidth]{fig15.png}
\caption{포메이션별·조건별 포획률(색: 지휘관, 빗금: 기동 계층, 오차막대: 시드 부트스트랩 신뢰구간)}
\label{fig:condbars}
\end{figure*}

Gemini 지휘관을 쓴 두 조건이 기준선 대비 유의한 개선을 보인다(각각 $+0.060$, $+0.062$).
반면 Qwen2.5 7B는 두 기동 조건 모두에서 기준선보다 \emph{유의하게 나쁘다}($-0.192$,
$-0.203$). 이 음의 결과는 \S\ref{sec:res_threshold}에서 다룬다.

\paragraph{포메이션별 분해}
\label{sec:res_ceiling}
통합 수치는 포메이션에 따라 다르게 읽어야 한다(Table~\ref{tab:ceiling}). 집중 포메이션은
휴리스틱만으로 포획률 0.997에 도달하므로 개선 방향으로는 어떤 방법도 변별되지 않고(기동 계층
단독 차이 $+0.000$, $95\,\%$ CI $[-0.010, +0.010]$), 통합값은 이 구간이 3분의 1을 차지하여
희석된다. 양동과 파상만 보면 제안 체계의 개선은 각각 $+0.090$, $+0.103$ 으로 통합값
$+0.062$ 보다 크다. 한편 이 구간은 악화 방향으로는 여전히 변별력을 가진다 --- Qwen2.5 7B
배정이 $-0.157$ ($[-0.323, -0.057]$), 같은 모델에 점수맵을 붙여도 $-0.110$
($[-0.247, -0.043]$), 14B에 점수맵을 붙인 경우가 $-0.040$ ($[-0.100, -0.010]$)으로, 약한
지휘관은 거의 완벽한 상황도 악화시킨다. 이는 \S\ref{sec:res_threshold}의 능력 임계가 가장
쉬운 포메이션에서도 관측된다는 뜻이다.

% 자동 생성 — tools/make_paper_tables.py. 손으로 고치지 말 것.
\begin{table}[t]
\centering
\caption{포메이션별 천장 효과. 집중 포메이션은 baseline 이 이미 0.997 이라 개선 여지가 0.003 뿐이다 --- 이 포메이션 단독으로 결론을 내면 안 된다.}
\label{tab:ceiling}
\footnotesize
\resizebox{\ifdim\width>\linewidth\linewidth\else\width\fi}{!}{%
\begin{tabular}{lcccc}
\toprule
포메이션 & baseline & 남은 여유 & 제안 & $\Delta$ \\
\midrule
집중 & 0.997 & 0.003 & 0.990 & -0.007 \\
양동 & 0.787 & 0.213 & 0.877 & +0.090$^{*}$ \\
파상 & 0.750 & 0.250 & 0.853 & +0.103 \\
\bottomrule
\end{tabular}
}
\\[3pt]
\begin{minipage}{\linewidth}\footnotesize
$^{*}$ 신뢰구간이 0 을 포함하지 않음.
\end{minipage}
\end{table}


\subsubsection{지휘관 능력 임계}
\label{sec:res_threshold}

기동 계층을 휴리스틱으로 고정하여 지휘관 축만 분리하면, 모델 능력에 따른 단조 증가가
나타난다(Fig.~\ref{fig:ladder}; 오차막대는 시드 부트스트랩 95\,\% 신뢰구간).

\begin{figure}[tbp]
\centering
\includegraphics[width=\columnwidth]{fig16.png}
\caption{지휘관 능력 사다리(기동 계층 휴리스틱 고정, 점선: 휴리스틱 배정 기준선)}
\label{fig:ladder}
\end{figure}

3개 포메이션 통합 기준으로 Qwen2.5 7B는 0.652, Qwen2.5 14B는 0.811, Gemini는 0.904이며,
휴리스틱 기준선은 0.844이다. 즉 \textbf{휴리스틱 선을 14B와 Gemini 사이에서 교차한다.}

이 결과가 함의하는 바는 ``LLM 지휘관을 쓰면 좋다''가 아니라 \textbf{``일정 능력 수준에 못
미치는 지휘관은 코드 휴리스틱보다 못하다''}이다. Qwen2.5 7B의 $-0.192$($\dz = -1.00$)는
유의한 악화이며, 이 음의 결과가 임계 주장을 신뢰할 수 있게 만든다.

\subsubsection{무엇이 지휘관 성능을 만드는가}
\label{sec:res_mechanism}

능력 순위 자체는 모델이 바뀌면 낡는다. 더 오래 남는 질문은 \emph{무엇이} 성능을 만드는가이다.
계획 단위 계측 결과를 Table~\ref{tab:planquality}와 Fig.~\ref{fig:planquality}에 정리하였다(점
하나가 지휘관 모델 하나, 오차막대는 양방향 부트스트랩 95\,\% 신뢰구간). 그림에 회귀선이나 상관계수는 얹지 않았다 --- 모델이 3종이라 $n=3$ 이며, 여기에 상관을
적으면 없는 통계적 근거를 주장하는 셈이 된다.

\begin{figure*}[!tbp]
\centering
\includegraphics[width=0.85\textwidth]{fig17.png}
\caption{계획 품질 대 성능: (a) 배정 커버리지, (b) 재계획 간 churn, (c) 경로 교차}
\label{fig:planquality}
\end{figure*}

% 자동 생성 — tools/make_paper_tables.py. 손으로 고치지 말 것.
\begin{table*}[t]
\centering
\caption{지휘관 계획 품질과 성능. 계획 품질은 시뮬레이션 결과와 독립적으로, 계획이 산출된 시점에 측정한다. \textbf{7B 는 커버리지가 가장 높은데 포획률은 가장 낮다} --- 실패 원인은 미배정이 아니다.}
\label{tab:planquality}
\footnotesize
\setlength{\tabcolsep}{4pt}%
\resizebox{\ifdim\width>\linewidth\linewidth\else\width\fi}{!}{%
\begin{tabular}{lcccccc}
\toprule
지휘관 & 지연 (s) & 커버리지 $\uparrow$ & churn $\downarrow$ & 교차 $\downarrow$ & 폴백 $\downarrow$ & 포획률 $\uparrow$ \\
\midrule
Qwen2.5 7B (로컬) & 2.75 & 0.934 & 0.324 & 0.152 & 0.007 & 0.652 \\
Qwen2.5 14B (로컬) & 16.59 & 0.877 & 0.167 & 0.111 & 0.001 & 0.811 \\
Gemini 3.5 Flash-Lite (클라우드) & 1.56 & 0.881 & 0.098 & 0.033 & 0.000 & 0.904 \\
\bottomrule
\end{tabular}
}
\\[3pt]
\begin{minipage}{\linewidth}\footnotesize
커버리지 $=$ 배정된 활성 클러스터 비율. churn $=$ 재계획 간 담당이 바뀐 배의 비율. 교차 $=$ 두 방어정의 요격 경로가 교차한 비율. 포획률은 기동 계층을 휴리스틱으로 고정한 조건(LLM 지휘관 + 휴리스틱 기동)의 값이다.
\end{minipage}
\end{table*}


결정적인 관측은 다음과 같다. \textbf{Qwen2.5 7B는 커버리지가 가장 높은데(0.934) 포획률은
가장 낮다(0.652).} 즉 막아야 할 무리를 배정하지 않아서 지는 것이 아니다. 오히려 7B는 거의
모든 무리에 배를 지목한다.

성능과 단조로 대응하는 것은 churn과 경로 교차이다. churn은 $0.324 \to 0.167 \to 0.098$,
경로 교차는 $0.152 \to 0.111 \to 0.033$ 으로 능력축을 따라 감소하며, 같은 순서로 포획률은
$0.652 \to 0.811 \to 0.904$ 로 증가한다(모두 기동 계층을 휴리스틱으로 고정한 조건). churn 0.324는 재계획마다 방어정의 약 3분의 1이
담당을 바꾼다는 뜻이며, 이는 \emph{반쯤 전개한 그물벽을 버리고 떠난다}는 것을 의미한다.
경로 교차는 두 방어정이 서로의 진로를 가로지르는 배정으로, 충돌 위험으로 직결된다.

이 관찰은 순위표가 아니라 설계 원리를 제시한다. \textbf{주기적으로 재계획하는 구조에서는
계획의 최적성보다 계획의 안정성이 중요하다.} 매 재계획마다 국소적으로 더 나은 배정을 찾아
바꾸는 것은, 그 전환 비용(버려지는 그물, 교차하는 경로)을 고려하면 손해일 수 있다.

\subsubsection{정성적 예시}
\label{sec:res_qualitative}

Fig.~\ref{fig:runshot}은 제안 체계가 실제로 동작하는 한 장면이다. 지휘관이 세 무리에 세
방어정을 각각 배정하고(패널 우측의 배정과 근거; 근거 문장은 모델 출력을 영문으로 옮겨
실었다), 각 방어정이 배정된 회랑에 그물벽을 친
결과 이 시점까지 적 3척이 포획되고 돌파와 아군 충돌은 없다. 재계획 주기는
\S\ref{sec:experiment}의 평가 설정과 같은 25\,step이다.

\textbf{운용 함의.} 지휘관 계층은 능력 임계 위의 모델에만 켠다 --- 본 실험 기준 로컬 7B는 끄는 편이 낫고 14B는 이득이 없으며, 클라우드 모델만 이득이었다. 지휘관을 고르거나 프롬프트를 손볼 때의 기준은 배정 커버리지가 아니라 재계획 간 담당 유지율이다.

\subsection{두 기여는 더해진다}
\label{sec:res_g3}

\subsubsection{두 계층의 기여는 가산적이다}
\label{sec:res_interaction}

식~\eqref{eq:interaction}의 상호작용을 10개 지표 전부에 대해 계산하였다
(Table~\ref{tab:interaction}, Fig.~\ref{fig:interaction}). 그림의 (a)에서 두 선이 평행에
가까울수록 두 계층의 기여가 가산적이고, (b)의 속이 빈 표식은 95\,\% 구간이 0을 포함하는
차이다. (a)에 오차막대를 두지 않은 것은 의도적이다 --- 셀별 구간은
대응이 풀린 구간이라 시나리오 난이도의 분산을 그대로 포함해 넓고, 그것을 겹쳐 그리면
유의한 대응차이가 유의하지 않은 것처럼 읽힌다. 불확실성은 전부 (b)에 있다.

% 자동 생성 — tools/make_paper_tables.py. 손으로 고치지 말 것.
\begin{table*}[tp]
\centering
\caption{2$\times$2 상호작용 분해 (Gemini 지휘관). 10개 지표 중 \textbf{9개}에서 상호작용 신뢰구간이 0 을 포함한다 --- 두 계층의 기여는 \textbf{가산적}이다.}
\label{tab:interaction}
\footnotesize
\resizebox{\ifdim\width>\linewidth\linewidth\else\width\fi}{!}{%
\begin{tabular}{lccc}
\toprule
지표 & 총효과 & 상호작용 [95\,\% CI] & 판정 \\
\midrule
포획률 & +0.062 & -0.039 [-0.083, +0.012] & 가산 \\
돌파 수 & -0.62 & +0.39 [-0.13, +0.83] & 가산 \\
충돌률(척당) & -0.052 & -0.015 [-0.070, +0.022] & 가산 \\
그물 걸림 & -0.60 & +0.17 [-0.08, +0.41] & 가산 \\
포획당 그물 & -0.328 & +0.004 [-0.045, +0.059] & 가산 \\
포획당 이동거리 (m) & -182 & -55 [-125, +30] & 가산 \\
총 선회량 (rad) & -757 & -56 [-165, +59] & 가산 \\
포획 이격거리 (m) & +223 & +38 [-2, +81] & 가산 \\
평균 포획 시각 (step) & -22 & -12 [-20, -6] & 상쇄 \\
전개 시 적 거리 (m) & -371 & +26 [-44, +109] & 가산 \\
\bottomrule
\end{tabular}
}
\\[3pt]
\begin{minipage}{\linewidth}\footnotesize
상호작용 $=$ (LLM+U-Net $-$ LLM+휴리스틱) $-$ (휴리스틱+U-Net $-$ baseline). 신뢰구간이 0 을 포함하면 두 계층이 서로를 돕지도 잡아먹지도 않는다.
\end{minipage}
\end{table*}


\begin{figure*}[tbp]
\centering
\includegraphics[width=0.78\textwidth]{fig18.png}
\caption{$2\times2$ 상호작용(Gemini 지휘관): (a) 네 셀의 포획률 평균, (b) 시드 대응차이와 BCa 신뢰구간}
\label{fig:interaction}
\end{figure*}

\textbf{10개 지표 중 9개에서 상호작용의 신뢰구간이 0을 포함한다.} 즉 두 계층은 서로를 돕지도,
서로를 잠식하지도 않는다. 본 논문은 이를 상승효과로 포장하지 않고 \emph{가산적}이라고
보고한다. 이는 약점이 아니라 설계 검증으로 읽어야 한다. 두 계층이 서로 다른 실패 모드를
다루므로 기여가 독립적으로 누적된다는 뜻이기 때문이다.

이 판정은 다중비교 보정에도 흔들리지 않는다. Bonferroni 보정 구간(99.5\,\%)으로 바꾸어도
0을 포함하는 지표는 그대로 9개다 --- 보정은 구간을 넓히므로 ``차이 없음''을 더 쉽게 만드는
방향인데, 그 방향으로 옮겨도 판정이 바뀌지 않는다는 뜻이다.

한 가지 예외는 평균 포획 시각이다. 상호작용이 $-11.98$ 이고 $95\,\%$ 신뢰구간이
$[-19.84,\; -5.93]$ 로 유의한 상쇄를 보인다. \textbf{이 예외는 보정 구간에서도
살아남는다} --- 99.5\,\% 구간이 $[-23.31,\; -3.54]$ 이다. 따라서 다중비교가 만들어
낸 우연이 아니다. 이는 해석이 된다. 두 계층 모두 제압을 앞당기는 방향으로 작용하는데,
\emph{같은 포획을 두 번 앞당길 수는 없다}. 공유 메커니즘에 대한 수확체감이다.

포메이션별로 보면 두 계층의 분업이 드러난다. 양동 포메이션에서는 기동 계층 단독
($-0.033$)도 전략 계층 단독($+0.070$)도 신뢰구간이 0을 포함하여 유의하지 않은데, 둘을
결합하면 $+0.090$ 으로 유의해진다. 각각으로는 부족한 두 기여가 합쳐져 임계를 넘는 것이며,
2계층 설계가 값을 하는 지점이 여기다. 이는 이 포메이션의 구조와 부합한다 --- 세 무리가
동시에 도래하고 방어정도 세 척이라 여유가 전혀 없으므로, 배정과 기동 어느 한쪽만 좋아서는
부족하고 둘 다 맞아야 막힌다. 반대로 파상 포메이션에서는 단이 시간차를 두고 오므로 자원
재배분의 여지가 있고, 기동 계층 단독($+0.157$)이 이미 유의하며 결합이 그보다 낫지 않다.

\subsubsection{요약: 무엇이 남는가 --- 포획률은 가장 둔한 지표다}
\label{sec:res_effectsize}

\S\ref{sec:metrics}에서 정의한 10개 지표 전부에 대해 제안 체계와 기준선의 대응 차이를
계산하였다(Table~\ref{tab:effectsize}, Fig.~\ref{fig:forest}; \texttt{llm\_unet}@Gemini 대
\texttt{heur\_heur}, 굵은 선은 95\,\%, 옅은 선은 Bonferroni 보정 99.5\,\% 구간, 속이 빈
표식은 보정 구간이 0을 포함). 그림의 $\dz$ 구간은 평균차의 BCa 신뢰구간을 $\mathrm{sd}(d)$ 로 나눈 것이므로, 그림이
0 을 배제하는 것과 표의 평균차 구간이 0 을 배제하는 것은 같은 사건이다.

% 자동 생성 — tools/make_paper_tables.py. 손으로 고치지 말 것.
\begin{table*}[t]
\centering
\caption{전 지표 효과크기 (제안 = LLM 지휘관 + U-Net 점수맵, Gemini). $|d_z|$ 내림차순. \textbf{포획률이 가장 둔한 지표이며, 다중비교 보정에서 유의성을 잃는 유일한 방어 성과 지표다.}}
\label{tab:effectsize}
\footnotesize
\resizebox{\ifdim\width>\linewidth\linewidth\else\width\fi}{!}{%
\begin{tabular}{lccccc}
\toprule
지표 & 방향 & baseline & 제안 & $\Delta$ & $d_z$ \\
\midrule
총 선회량 (rad) & $\downarrow$ & 1625 & 868 & -757$^{*}$$^{\dagger}$ & -3.23 \\
포획당 그물 & $\downarrow$ & 0.737 & 0.409 & -0.328$^{*}$$^{\dagger}$ & -2.09 \\
포획 이격거리 (m) & $\uparrow$ & 1511 & 1734 & +223$^{*}$$^{\dagger}$ & +1.81 \\
전개 시 적 거리 (m) & $\downarrow$ & 1499 & 1128 & -371$^{*}$$^{\dagger}$ & -1.60 \\
그물 걸림 & $\downarrow$ & 0.82 & 0.22 & -0.60$^{*}$$^{\dagger}$ & -1.19 \\
평균 포획 시각 (step) & $\downarrow$ & 430 & 408 & -22$^{*}$$^{\dagger}$ & -1.04 \\
포획당 이동거리 (m) & $\downarrow$ & 825 & 643 & -182$^{*}$$^{\dagger}$ & -0.65 \\
충돌률(척당) & $\downarrow$ & 0.130 & 0.078 & -0.052$^{*}$$^{\dagger}$ & -0.52 \\
돌파 수 & $\downarrow$ & 1.56 & 0.93 & -0.62$^{*}$ & -0.44 \\
포획률 & $\uparrow$ & 0.844 & 0.907 & +0.062$^{*}$ & +0.44 \\
\bottomrule
\end{tabular}
}
\\[3pt]
\begin{minipage}{\linewidth}\footnotesize
방향 $\uparrow$ = 높을수록 좋음, $\downarrow$ = 낮을수록 좋음. $^{*}$ 95\,\% 신뢰구간이 0 을 포함하지 않음. $^{\dagger}$ 지표 10개에 대한 Bonferroni 보정(99.5\,\% 구간)에서도 0 을 배제함 --- 즉 단검이 없는 행은 보정을 견디지 못한다.
\end{minipage}
\end{table*}


\begin{figure}[!tp]
\centering
\includegraphics[width=\columnwidth]{fig19.png}
\caption{제안 대 기준선의 전 지표 효과크기 $\dz$(굵은 선 신뢰구간, 옅은 선 Bonferroni 보정 구간)}
\label{fig:forest}
\end{figure}

결과는 명확하다. \textbf{포획률의 효과크기($\dz = 0.44$)는 측정한 10개 지표 중 최하위권이며,
총 선회량($-3.23$)의 7분의 1 수준이다.} 개선이 가장 큰 지표는 기동 비용과 자원 효율 쪽이다.

\begin{itemize}[leftmargin=*, itemsep=2pt, topsep=3pt]
\item 총 선회량 $1625 \to 868$ ($-47\,\%$, $\dz = -3.23$)
\item 포획당 그물 $0.737 \to 0.409$ ($-45\,\%$, $\dz = -2.09$)
\item 포획 이격거리 $1511 \to 1734\unit{m}$ ($\dz = +1.81$)
\item 그물 걸림 $0.82 \to 0.22$ ($-73\,\%$, $\dz = -1.19$)
\end{itemize}

\textbf{이 대조는 다중비교 보정에서 더 선명해진다.} 지표 10개에 대한 Bonferroni 보정
(99.5\,\% 구간)을 적용하면 10개 중 8개가 그대로 유의하지만, \textbf{포획률과 돌파 수는
유의성을 잃는다}(Table~\ref{tab:effectsize}의 단검 표시). 포획률의 보정 구간은
$[-0.003, +0.132]$ 로 0을 아슬아슬하게 포함한다. 기동 계층 단독 조건에서도 같은 두 지표가
탈락한다($[-0.013, +0.100]$).

이 결과를 감출 이유가 없다. 오히려 이 논문이 말하려는 바를 정확히 드러내기 때문이다 ---
\emph{보정을 견디지 못하는 것은 ``더 많이 막는다''는 주장이고, 견디는 것은 ``같은 정도로
막으면서 자원과 기동과 위험을 크게 줄인다''는 주장이다.} 따라서 본 논문은 포획률 개선을
주장의 중심에 두지 않는다.

즉 제안 체계의 이득은 ``더 많이 잡는다''보다 \textbf{``같은 정도로 막으면서 훨씬 적은 자원과
기동, 더 낮은 위험으로 해낸다''}에 가깝다. 이는 자원이 제한된 실운용에서 오히려 더 중요한
성질이다. 포획률의 효과크기가 작은 것은 방법의 한계가 아니라 지표의 천장 때문이며
(\S\ref{sec:res_ceiling}), 효율 지표에는 그러한 천장이 없다.

\textbf{운용 함의.} 두 계층이 서로 간섭하지 않으므로 각 계층을 독립적으로 개선·교체할 수 있다. 성과 관리 지표로는 천장이 있는 포획률보다 자원·기동 지표(포획당 그물, 총 선회량)가 변별력이 있다.

\subsection{배치 관점: 지연}
\label{sec:res_latency}

응답 지연은 모델 규모에 단조가 아니다(Table~\ref{tab:planquality}). 로컬 14B가 16.59초인 데
비해 클라우드 모델은 1.56초로, \textbf{로컬이 10.6배 느리다}. 한 결정 주기가 25초
(25\,step $\times$ $\Delta t\!=\!1\unit{s}$)이므로 14B의 평균 지연은 주기의 \textbf{약
3분의 2}를 소모한다.

평균만 보면 주기 안에 들어오지만, 결정적인 것은 분포의 꼬리다
(Table~\ref{tab:latency_tail}, Fig.~\ref{fig:latency_ecdf}). 14B 호출 5{,}366건 가운데 \textbf{4.2\,\%가 25초 주기를
초과}하였고, 99 백분위는 27.8초, 최댓값은 296.4초였다. 동기 구조라면 이 4.2\,\%마다 제어
루프가 실제로 정지한다. 반면 클라우드 모델은 최댓값이 2.75초로 주기의 11\,\%에 그쳐
동기 호출로도 여유가 있다.

\begin{figure}[!tbp]
\centering
\includegraphics[width=\columnwidth]{fig20.png}
\caption{지휘관 응답 지연의 경험 누적분포(파선은 결정 주기)}
\label{fig:latency_ecdf}
\end{figure}

\begin{table*}[tp]
\centering
\caption{지휘관 응답 지연 분포. 결정 주기 25초와 비교한다. 로컬 14B만 주기를 초과하는
         꼬리를 갖는다.}
\label{tab:latency_tail}
\footnotesize
\begin{tabular}{lrrrrr}
\toprule
지휘관 & $n$ & 평균 & $p_{99}$ & 최대 & $>25$초 \\
\midrule
Qwen2.5 7B (로컬)    & 4{,}299 & 2.75  & 7.08  & 10.06  & 0.0\,\% \\
Qwen2.5 14B (로컬)   & 5{,}366 & 16.59 & 27.81 & 296.41 & \textbf{4.2\,\%} \\
Gemini (클라우드)    & 3{,}882 & 1.56  & 2.30  & 2.75   & 0.0\,\% \\
\bottomrule
\end{tabular}
\\[3pt]
\begin{minipage}{\linewidth}\footnotesize
단위는 초. 로컬 모델은 Ollama 서빙(Qwen2.5 14B는 Q4\_K\_M 양자화, 문맥 16{,}384 토큰),
NVIDIA RTX 4070(12\,GB) 단일 GPU에서 측정하였다.
\textbf{★ 해석 시 주의}: 이 GPU에는 시뮬레이터와 정책 추론이 동시에 올라가 있었다.
14B는 Q4 가중치만 약 9\,GB이므로 KV 캐시를 더하면 12\,GB를 압박하며, 최댓값 296초는
가중치 오프로드/재적재를 시사한다. 실제로 본 저장소의 설치 문서는 같은 GPU에서
$2\sim4$초를 예상하고 있어 실측과 $4\sim8$배 차이가 난다. 따라서
\textbf{14B의 꼬리가 로컬 중형 모델 서빙의 본질적 성질인지, 이 구성에서 GPU를 공유한
결과인지는 본 실험으로 구분되지 않는다}. 여기서 읽어야 할 것은 개별 수치가 아니라
\emph{자원이 제한된 온보드 환경에서는 로컬 중형 모델의 지연이 결정 주기를 넘길 수 있다}는
정성적 사실이며, 이는 전용 가속기를 갖추면 완화될 수 있다.
\end{minipage}
\end{table*}

단 이 수치의 측정 구성에는 교란 요인이 있다(\S\ref{sec:limitations} 참조). 즉 \S\ref{sec:arch}의 비동기 설계는 클라우드 배치에서는 선택이지만 \textbf{로컬 배치에서는
지연 꼬리에 대한 안전장치}로서 실질적 의미를 갖는다.

\textbf{운용 함의.} 로컬 배치에서 비동기 분리는 선택이 아니라 전제이며, 배치 전에 확인할 값은 평균 지연이 아니라 주기를 넘기는 호출의 비율이다.
